\chapter{Analisi as-is: il dataset di partenza}
\label{chap:dataset}

\paragraph{Riassunto del capitolo.} Questo capitolo è l'\emph{analisi
as-is}: descrive con precisione i dati così come ci sono stati
consegnati dal Politecnico. Il file MATLAB grezzo viene convertito in
un CSV \emph{tidy} di \num{9805} righe e 6 colonne (Aging, Temperatura,
SOC, Frequenza, Re/Im di $Z$). Quattro invarianti sono verificati da
test automatici. L'esplorazione preliminare mostra una topologia di
Nyquist coerente con la teoria e una separabilità young/old già
visibile a occhio: due osservazioni che giustificano \emph{ex ante} la
fattibilità dei tre task definiti nel \cref{chap:introduction}.

\section{Origine e formato grezzo}

Il dato grezzo arriva al progetto come file MATLAB
(\codefile{data/raw/GEIS.mat}) prodotto dallo strumento Bio-Logic. La
struttura interna è gerarchica:

\begin{lstlisting}[language=Matlab]
GEIS.mat
 |-- Aging0
 |    |-- GEIS_20    (matrice [n_freq x 5 SOC x 3])
 |    |-- GEIS_22_5  ...
 |    |-- ...
 |    |-- GEIS_47_5
 |-- Aging1
 |-- ...
 |-- Aging4
\end{lstlisting}

Ogni matrice 3D contiene, per ciascuno dei 5 livelli di SOC, le triplette
$(\text{Frequency [Hz]}, \text{Re(Z) [\(\Omega\)]}, \text{Im(Z) [\(\Omega\)]})$.
I valori grezzi sono in $\Omega$; noi moltiplichiamo per $10^3$ per
ottenere milliohm, una scelta più leggibile per i plot di Nyquist su
celle da pochi \milliohm{} di resistenza serie.

\section{Pipeline di pulizia}

La conversione è identica nel notebook
\codefile{notebooks/0\_mat\_to\_csv.ipynb} (variante interattiva,
pedagogica) e in \codefile{src/data/preprocessing.py::mat\_to\_dataframe}
(variante di produzione, testata). I passaggi in entrambi sono:
\begin{enumerate}[leftmargin=*]
  \item caricamento con \texttt{scipy.io.loadmat(squeeze\_me=True,
        struct\_as\_record=False)};
  \item iterazione su \texttt{Aging\{0..4\}.GEIS\_<temperature>};
  \item per ciascun SOC, espansione delle righe in record
        $(Aging, Temperature, SOC, Frequency, Z_\text{real}, Z_\text{imag})$;
  \item conversione $\Omega \to \milliohm{}$ tramite fattore $10^3$;
  \item salvataggio CSV \emph{tidy} in
        \codefile{data/processed/batteries\_cleaned\_dataset.csv}.
\end{enumerate}

Il risultato è un DataFrame da \num{9805} righe con sei colonne, riassunto
in \cref{tab:dataset-schema}.

\begin{table}[H]
  \centering
  \caption{Schema del dataset processato.}
  \label{tab:dataset-schema}
  \begin{tabular}{lll}
    \toprule
    \textbf{Colonna} & \textbf{Tipo} & \textbf{Range / unità} \\
    \midrule
    \texttt{Aging}        & int   & $\{0, 1, 2, 3, 4\}$ (fresh $\to$ aged) \\
    \texttt{Temperature}  & float & 8 livelli in \si{\celsius}: 20.0, 22.5, 25.0, 27.5, 30.0, 35.0, 40.0, 47.5 \\
    \texttt{SOC}          & int   & $\{0, 1, 2, 3, 4\}$ ($\approx$ 0\,\%, 25\,\%, 50\,\%, 75\,\%, 100\,\%) \\
    \texttt{Frequency}    & float & \si{\hertz}, $\sim$49 valori log-spaziati in $[0.099, 10\,002]$ \\
    \texttt{Z\_real}      & float & $\mathrm{Re}(Z)$ in \milliohm{} \\
    \texttt{Z\_imag}      & float & $\mathrm{Im}(Z)$ in \milliohm{} \\
    \bottomrule
  \end{tabular}
\end{table}

\section{Sanity check}

Quattro invarianti vengono verificati dai test in
\codefile{tests/test\_data.py} e \codefile{tests/test\_preprocessing.py}:

\begin{enumerate}[leftmargin=*]
  \item $\texttt{df.shape} = (9805, 6)$;
  \item nessun valore mancante: $\texttt{df.isna().sum().sum()} = 0$;
  \item esattamente 200 gruppi
        $(Aging, Temperature, SOC)$ -- 40 coppie ($A$, $T$) per 5 SOC;
  \item $\texttt{Frequency}$ strettamente positiva e contenuta in
        $[0.099, 10\,002]\,\si{\hertz}$.
\end{enumerate}

\section{Esplorazione preliminare}

\begin{figure}[htbp]
  \centering
  \includegraphics[width=0.68\linewidth]{dataset_nyquist.pdf}
  \caption{Curve di Nyquist al livello di invecchiamento centrale
  ($Aging = 2$, $\text{SOC} = 2$) per le otto temperature del dataset.
  La forma a semicerchio + coda diffusiva è la firma elettrochimica
  della cella; cresce in ampiezza al diminuire della temperatura.}
  \label{fig:dataset-nyquist}
\end{figure}

\begin{figure}[htbp]
  \centering
  \includegraphics[width=0.68\linewidth]{dataset_aging_evolution.pdf}
  \caption{Evoluzione delle curve di Nyquist con il livello di
  invecchiamento, a $T = \SI{25}{\celsius}$ e SOC $=2$. Il semicerchio
  si allarga e la coda si trasla a destra al crescere dell'Aging --
  giustificazione empirica della separabilità che il Task 2 sfrutterà.}
  \label{fig:dataset-aging-evolution}
\end{figure}

\begin{figure}[htbp]
  \centering
  \includegraphics[width=0.68\linewidth]{dataset_boxplot_zreal.pdf}
  \caption{Distribuzione di $\mathrm{Re}(Z)$ stratificata per livello
  di Aging, aggregata su tutte le temperature, SOC e frequenze. La
  mediana sale monotonamente con l'invecchiamento.}
  \label{fig:dataset-boxplot}
\end{figure}

Le prime osservazioni guidano le scelte di feature engineering del
prossimo capitolo:
\begin{itemize}
  \item Le curve di Nyquist a parità di Aging hanno topologia simile
        (semicerchio + coda) ma \emph{traslate} verso destra al crescere
        dell'aging, conferma del fatto che $R_\mathrm{s}$ e
        $R_\mathrm{ct}$ aumentano con l'invecchiamento.
  \item La media di $\mathrm{Re}(Z)$ in funzione di $T$ ha un andamento
        monotono decrescente in alta frequenza (cinetica più rapida) e
        meno netto in bassa frequenza, dove dominano i processi
        diffusivi.
  \item A SOC fissato, le distribuzioni di $Z_\text{real}$ per cellule
        \emph{young} (Aging 0--2) sono visibilmente separate da quelle
        \emph{old} (Aging 3--4) -- giustificazione empirica per la natura
        binaria del Task~2.
\end{itemize}

\section{Caching e accesso}

\codefile{src/data/loader.py::load\_dataset} è decorato con
\texttt{@functools.lru\_cache(maxsize=1)}: il CSV viene letto una sola
volta per processo. Il loader ordina sempre per $(Aging, Temperature,
SOC, Frequency)$ in modo che split successivi (anche complessi come il
\acrshort{logo}) restituiscano indici stabili e riproducibili.
